% Standalone source. Compile twice with: pdflatex mi28_full_proof.tex
\documentclass[11pt,a4paper]{article}
\usepackage[T1]{fontenc}
\usepackage[utf8]{inputenc}
\usepackage{lmodern}
\usepackage[margin=27mm]{geometry}
\usepackage{amsmath,amssymb,amsthm,mathtools}
\usepackage{microtype}
\usepackage{enumitem}
\usepackage{needspace}
\usepackage{xurl}
\usepackage[hidelinks]{hyperref}
\hypersetup{
  pdftitle={MI-28: A determinant inequality for all nonnegative base powers},
  pdfauthor={George Stepaniants},
  pdfsubject={A log-majorization proof using Furuta inequalities and parameter swapping}
}
\setlength{\parskip}{0.3em}
\setlength{\emergencystretch}{2em}
\setlist[enumerate]{itemsep=0.35em,topsep=0.4em}
\newtheorem{theorem}{Theorem}
\newtheorem{lemma}[theorem]{Lemma}
\newtheorem{corollary}[theorem]{Corollary}
\theoremstyle{remark}
\newtheorem{remark}[theorem]{Remark}
\newcommand{\norm}[1]{\left\lVert #1\right\rVert}
\newcommand{\lm}{\prec_{\log}}
\newcommand{\C}{\mathbb C}

\title{\textbf{MI-28: A determinant inequality\\
for all nonnegative base powers}\\[0.4em]
\large Furuta inequalities and an interchange of the parameters}
\author{George Stepaniants\\[0.3em]
\small Department of Computing and Mathematical Sciences\\
\small California Institute of Technology\\
\small Pasadena, California, USA\\
\small \href{mailto:gstepan@caltech.edu}{\texttt{gstepan@caltech.edu}}}
\date{11 September 2026}

\begin{document}
\maketitle

\begin{abstract}
For positive definite matrices \(A,B\), every \(k\ge0\), and
\(0\le p\le2\), we prove
\[
 \det(A^k+|AB|^p)\ge\det(A^k+A^pB^p).
\]
In fact, a log-majorization of the corresponding normalized matrices
holds. The previously proved range \(k\ge2\) is due to
Ghabries, Abbas, Mourad, and Assi. We treat the remaining base powers
using a normalized order implication. Two applications of the Furuta
inequality establish complementary parameter regions. The identity
\(\lvert |AB|^{-1}B\rvert=A^{-1}\) interchanges the two parameters
and closes the remaining range.
\end{abstract}

\noindent\textbf{Provenance and review.}
This manuscript was prepared with assistance from OpenAI Codex.
The accompanying review records identify the exact version checked
by independent Codex agents. Such review is neither human peer review
nor formal verification.

\section{Statement and notation}

Throughout, \(A,B\in\C^{n\times n}\) are positive definite,
matrix powers are defined by spectral functional calculus, and
\(\norm{\cdot}\) denotes the operator norm.
For a matrix \(M\), write \(|M|=(M^*M)^{1/2}\).
Eigenvalues of a positive definite matrix are listed in decreasing order.
For positive vectors \(u,v\in(0,\infty)^n\), the notation
\(u\lm v\) means
\[
 \prod_{i=1}^j u_i\le\prod_{i=1}^j v_i\quad(1\le j<n),
 \qquad
 \prod_{i=1}^n u_i=\prod_{i=1}^n v_i .
\]

\begin{theorem}\label{thm:main}
For every \(A,B>0\), \(k\ge0\), and \(0\le p\le2\),
\begin{equation}\label{eq:mainlog}
 \lambda\!\left(A^{(p-k)/2}B^pA^{(p-k)/2}\right)
 \lm
 \lambda\!\left(A^{-k/2}|AB|^pA^{-k/2}\right).
\end{equation}
Consequently,
\begin{equation}\label{eq:maindet}
 \det(A^k+|AB|^p)\ge\det(A^k+A^pB^p).
\end{equation}
\end{theorem}

Inequality~\eqref{eq:maindet} is the positive definite formulation
of MI-28~\cite{MI28,GThesis}. Although \(A^k+A^pB^p\) need not be
Hermitian, its determinant is positive: after factoring out \(A^k\),
the remaining product \(A^{p-k}B^p\) is similar to the positive
definite matrix on the left of~\eqref{eq:mainlog}.
The range \(k\ge2\), \(0\le p\le2\) follows already from
\cite[Lemma~2.5]{GAMA20}; we record the precise substitution in
Section~\ref{sec:finish}.

\section{Two order implications}

We use the L\"owner--Heinz inequality
\begin{equation}\label{eq:LH}
 0<X\le Y,\quad 0\le t\le1
 \quad\Longrightarrow\quad X^t\le Y^t,
\end{equation}
and the following convention for the Furuta inequality~\cite{Furuta87}:
if \(X\ge Y>0\), \(r\ge0\), \(a\ge0\), \(q\ge1\), and
\((1+r)q\ge a+r\), then
\begin{equation}\label{eq:Furuta}
 (X^{r/2}Y^aX^{r/2})^{1/q}\le X^{(a+r)/q}.
\end{equation}
Both facts are also restated in
\cite[Propositions~1--2]{Tanahashi99}. The \(r\) used here is twice
the sandwich exponent in that reference.

For \(k,p>0\), let \(\mathcal I(k,p)\) denote the assertion that,
for every \(A,B>0\), with \(D=|AB|\),
\begin{equation}\label{eq:imp}
 D^p\le A^k\quad\Longrightarrow\quad B^p\le A^{k-p}.
\end{equation}
This notation describes a universal order implication.
It makes no assumption that \(A\) and \(B\) commute.

\begin{lemma}\label{lem:small}
If \(k>0\) and \(k/(k+1)\le p\le1\), then \(\mathcal I(k,p)\)
holds.
\end{lemma}

\begin{proof}
Assume \(D^p\le A^k\). In~\eqref{eq:Furuta}, choose
\[
 X=A^k,\qquad Y=D^p,\qquad r=\frac2k,\qquad
 a=\frac2p,\qquad q=2 .
\]
The condition \((1+r)q\ge a+r\) is exactly
\(p\ge k/(k+1)\). Thus
\[
 (AD^2A)^{1/2}\le A^{1+k/p}.
\]
Since \(D^2=BA^2B\), we have \(AD^2A=(ABA)^2\), and its positive
square root is \(ABA\). Congruence by \(A^{-1}\) gives
\(B\le A^{k/p-1}\). Applying~\eqref{eq:LH} with exponent
\(p\in(0,1]\) proves \(B^p\le A^{k-p}\).
\end{proof}

\begin{lemma}\label{lem:large}
If \(k>0\) and \(1\le p\le2\), then \(\mathcal I(k,p)\) holds.
\end{lemma}

\begin{proof}
Assume \(D^p\le A^k\), and set
\[
 C=ABA,\qquad s=\frac{p(k+2)}{k+p}.
\]
The hypotheses give \(1\le p\le s\le2\).
Apply~\eqref{eq:Furuta} with
\[
 X=A^k,\qquad Y=D^p,\qquad r=\frac2k,\qquad
 a=\frac2p,\qquad q=\frac2s.
\]
Here \(q\ge1\) and \((1+r)q=a+r\). Using \(AD^2A=C^2\), we obtain
\begin{equation}\label{eq:Cbound}
 C^s=(AD^2A)^{s/2}\le A^{k+2}.
\end{equation}
Because \(0<2/(k+2)\le1\), L\"owner--Heinz followed by inversion
gives
\begin{equation}\label{eq:Aminus}
 A^{-2}\le C^{-2s/(k+2)}=C^{-2p/(k+p)}.
\end{equation}

For every invertible \(S\) and real \(p\), a singular-value
decomposition gives
\[
 (SS^*)^p=S(S^*S)^{p-1}S^*.
\]
Apply this identity to \(S=A^{-1}C^{1/2}\), for which \(SS^*=B\):
\[
 B^p=A^{-1}C^{1/2}
       (C^{1/2}A^{-2}C^{1/2})^{p-1}
       C^{1/2}A^{-1}.
\]
Congruence in~\eqref{eq:Aminus}, followed by
L\"owner--Heinz with \(0\le p-1\le1\), therefore yields
\begin{align}
 B^p
 &\le A^{-1}C^{\,1+(p-1)(k-p)/(k+p)}A^{-1}\notag\\
 &=A^{-1}C^{\,p(k-p+2)/(k+p)}A^{-1}.
 \label{eq:Bbound}
\end{align}
The exponent of \(C\) in~\eqref{eq:Bbound} equals \(s\theta\), where
\[
 \theta=\frac{k-p+2}{k+2}\in[0,1].
\]
Applying L\"owner--Heinz to~\eqref{eq:Cbound} gives
\(C^{s\theta}\le A^{k-p+2}\). Substitute this in~\eqref{eq:Bbound}
to conclude \(B^p\le A^{k-p}\).
\end{proof}

\begin{remark}
The argument after~\eqref{eq:Cbound} is a direct proof of the
needed specialization of the negative-power Furuta inequality.
In the convention of~\eqref{eq:Furuta}, the relevant parameters in
\cite[Theorem~3]{Tanahashi99} are
\[
 X=A^{k+2},\quad Y=C^s,\quad
 a=\frac{k+p}{p(k+2)},\quad q=\frac1p,\quad r=-\frac2{k+2}.
\]
The displayed proof above verifies that specialization using
only L\"owner--Heinz and the singular-value identity.
\end{remark}

\section{Interchanging the parameters}

\begin{lemma}\label{lem:swap}
Let \(0<p\le k\). If \(\mathcal I(p,k)\) holds, then
\(\mathcal I(k,p)\) holds.
\end{lemma}

\begin{proof}
Assume \(D^p\le A^k\), where \(D=|AB|\), and define
\[
 \widetilde A=D^{-1},\qquad \widetilde B=B.
\]
The order of the factors is essential. Directly,
\[
 |\widetilde A\widetilde B|^2
   =BD^{-2}B
   =B(B^{-1}A^{-2}B^{-1})B
   =A^{-2},
\]
so \(|\widetilde A\widetilde B|=A^{-1}\).
Inverting the assumed order gives
\[
 |\widetilde A\widetilde B|^k=A^{-k}
 \le D^{-p}=\widetilde A^p.
\]
Apply \(\mathcal I(p,k)\) to this pair:
\begin{equation}\label{eq:swapped}
 B^k\le \widetilde A^{p-k}=D^{k-p}.
\end{equation}
Since \(0<p/k\le1\), L\"owner--Heinz applied
to~\eqref{eq:swapped} gives
\[
 B^p\le D^{p(k-p)/k}.
\]
Since \(0\le(k-p)/k\le1\), applying it also to \(D^p\le A^k\)
gives
\[
 D^{p(k-p)/k}\le A^{k-p}.
\]
Together these inequalities prove \(\mathcal I(k,p)\).
If \(p=k\), the second exponent is zero and its inequality is
simply \(I\le I\); the argument still applies.
\end{proof}

\begin{corollary}\label{cor:allsmallk}
For \(0<k\le2\) and \(0<p\le2\), the implication
\(\mathcal I(k,p)\) holds.
\end{corollary}

\begin{proof}
If \(p\ge1\), use Lemma~\ref{lem:large}.
Suppose \(0<p<1\). If \(p\ge k/(k+1)\), use
Lemma~\ref{lem:small}. In the remaining case,
\[
 p<\frac{k}{k+1}<k\le2.
\]
We establish the swapped implication \(\mathcal I(p,k)\).
When \(k\le1\), its exponent \(k\) satisfies
\[
 \frac{p}{p+1}<p<k\le1,
\]
so Lemma~\ref{lem:small}, with the parameters interchanged,
applies. When \(1\le k\le2\), Lemma~\ref{lem:large}, with the
parameters interchanged, applies. Lemma~\ref{lem:swap} now proves
\(\mathcal I(k,p)\).
\end{proof}

\section{Log-majorization and the determinant}\label{sec:finish}

\begin{proof}[Proof of Theorem~\ref{thm:main}]
First take \(0<k\le2\), \(0<p\le2\), and write
\[
 H=A^{(p-k)/2}B^pA^{(p-k)/2},\qquad
 Z=A^{-k/2}D^pA^{-k/2}.
\]
Both matrices are homogeneous of degree \(p\) in \(B\).
Let \(c=\norm{Z}>0\), and replace \(B\) by \(c^{-1/p}B\).
The resulting \(Z\) is at most \(I\), equivalently
\(D^p\le A^k\) for the scaled pair.
Corollary~\ref{cor:allsmallk} gives \(B^p\le A^{k-p}\).
Congruence by \(A^{(p-k)/2}\) shows that the scaled \(H\) is
at most \(I\). Scaling back proves
\begin{equation}\label{eq:norm}
 \norm{H}\le\norm{Z}.
\end{equation}

For every \(1\le j\le n\), apply~\eqref{eq:norm} to
\(\bigwedge^j A\) and \(\bigwedge^j B\).
These exterior powers are positive definite and preserve products,
inverses, and spectral powers. In particular,
\[
 \left|(\bigwedge\nolimits^j A)(\bigwedge\nolimits^j B)\right|
 =\bigwedge\nolimits^j |AB|.
\]
The operator norm of an exterior power of a positive definite matrix
is the product of its \(j\) largest eigenvalues. We thus obtain
all partial-product inequalities in~\eqref{eq:mainlog}.
The full determinants agree:
\begin{equation}\label{eq:equaldets}
 \det H=\det Z=\det(A)^{p-k}\det(B)^p.
\end{equation}
This proves~\eqref{eq:mainlog} for \(0<k\le2\), \(0<p\le2\).

For \(k\ge2\), use the previously established result
\cite[Lemma~2.5]{GAMA20}. In its notation, for \(0\le t\le s\le K\),
\[
 \lambda(X^{K-t}Y^t)
 \prec_{\mathrm{wlog}}
 \lambda\!\left(
 X^{K/2}(Y^{s/2}X^{-s}Y^{s/2})^{t/s}X^{K/2}\right).
\]
Take \(X=A^{-1}\), \(Y=B\), \(K=k\), \(s=2\), and \(t=p\).
The two eigenvalue lists are exactly those of \(H\) and \(Z\);
the product \(A^{p-k}B^p\) is similar to \(H\).
Equation~\eqref{eq:equaldets} changes this weak log-majorization
into the log-majorization~\eqref{eq:mainlog}.
This invokes the original range \(0\le t\le s\le K\) without
altering any of its parameter restrictions.

When \(p=0\), the two matrices in~\eqref{eq:mainlog} are both
\(A^{-k}\), so equality holds.
For \(k=0\), \(p>0\), let \(k\downarrow0\) in the established
range \(0<k\le2\). Functional calculus and the ordered eigenvalues
are continuous on the positive definite cone, and each
partial-product inequality and~\eqref{eq:equaldets} persists.
Thus~\eqref{eq:mainlog} holds at every stated endpoint.

Finally, log-majorization is ordinary majorization of the logarithms
of the eigenvalues. Applying the convex function
\(t\mapsto\log(1+e^t)\) gives
\[
 \det(I+H)\le\det(I+Z).
\]
Indeed, the corresponding inequality for the sums of this convex
function is the standard convex-function characterization of
majorization. Since \(A^{p-k}B^p\) is similar to \(H\),
\begin{align*}
 \det(A^k+A^pB^p)&=\det(A^k)\det(I+H),\\
 \det(A^k+D^p)&=\det(A^k)\det(I+Z).
\end{align*}
Multiplying by the positive number \(\det(A^k)\) proves
\eqref{eq:maindet}.
\end{proof}

\begin{remark}[Scope]
The theorem treats every dimension and every parameter pair in
the canonical positive definite statement of MI-28.
For \(k>0\) and \(p>0\), the positive semidefinite extension follows
by applying the determinant inequality to \(A+\varepsilon I\) and
\(B+\varepsilon I\), then letting \(\varepsilon\downarrow0\).
The positive definite formulation avoids conventions for zeroth
powers of singular matrices.
\end{remark}

\Needspace{18\baselineskip}
\begin{thebibliography}{9}

\bibitem{MI28}
\texttt{MColbrook/OpenProblemsInNLA}.
\emph{MI-28: Ghabries's determinant comparison for arbitrary
nonnegative base powers}.
\url{https://github.com/MColbrook/OpenProblemsInNLA/blob/main/matrix-inequalities-and-norms/MI-28/README.md}.

\bibitem{GThesis}
M.~M.~Ghabries.
\emph{Contributions to Matrix Inequalities and Some Applications}.
PhD thesis, University of Angers and Lebanese University, 2022.
Section~2.5 and final Open Problems, Problem~1, pp.~109--110.
\url{https://www.researchgate.net/publication/361793582_Contributions_to_Matrix_Inequalities_and_Some_Applications}.

\bibitem{GAMA20}
M.~M.~Ghabries, H.~Abbas, B.~Mourad, and A.~Assi.
\emph{A proof of a conjectured determinantal inequality}.
Linear Algebra and its Applications \textbf{605} (2020), 21--28.
Lemma~2.5 and Theorem~1.1.
\url{https://doi.org/10.1016/j.laa.2020.07.013}.
Author-uploaded full text:
\url{https://www.researchgate.net/publication/342908148_A_proof_of_a_conjectured_determinantal_inequality}.

\bibitem{Furuta87}
T.~Furuta.
\emph{\(A\ge B\ge O\) assures
\((B^rA^pB^r)^{1/q}\ge B^{(p+2r)/q}\)
for \(r\ge0\), \(p\ge0\), \(q\ge1\) with
\((1+2r)q\ge p+2r\)}.
Proceedings of the American Mathematical Society
\textbf{101} (1987), 85--88.

\bibitem{Tanahashi99}
K.~Tanahashi.
\emph{The Furuta inequality with negative powers}.
Proceedings of the American Mathematical Society
\textbf{127} (1999), 1683--1692.
Propositions~1--2 and Theorem~3.
\url{https://doi.org/10.1090/S0002-9939-99-04705-X}.
Author-uploaded full text:
\url{https://www.researchgate.net/publication/255605812_The_Furuta_inequality_with_negative_powers}.

\end{thebibliography}
\end{document}
